%!TEX root = Vorlage_Buch.tex
\chapter{Leckeres aus dem Holzbackofen}\label{Chapter4}
\lettrine[lines=3]{D}{er Holzbackofen} ist eine echte Bereicherung für das 
Outdoor-Kochvergnügen.
Die Einsatzmöglichkeiten sind schier unerschöpflich. Fast alles das auf
dem Herd oder dem Grill gemacht wird, ist auch im Backofen möglich. Ob
Fleisch, Fisch, Gemüse, Pizzen, Brote oder Desserts das alles gelingt im 
Holzbackofen. Der Kreativität wird keine Grenzen gesetzt. Das Charmante an
dieser Zubereitungsmethode ist die lange Haltezeit der Temperatur. Pizzen
werden jenseits der 300 ℃ gebacken, danach kann locker noch ein 
Holzofenbrot gebacken werden. So kann die Restwärme noch sinnvoll genutzt
werden.

\begin{figure}[htbp]
	\centering
	\begin{minipage}{1\textwidth}
		\centering
		\includegraphics[width=.9\linewidth]{pics/Backofen}
		\captionof{figure}{HolzBackofen}
		\label{fig:Backofen}
	\end{minipage}
\end{figure}
\newpage

\section{Teige und was sich daraus machen lässt}

\subsection{Eine echt leckere Pizza}

\paragraph{Geräte}

\begin{itemize}[noitemsep]
	\item Holzbackofen
	\item Kenwood Cooking Chef
\end{itemize}

\paragraph{Zutaten Teig}

\begin{itemize}[noitemsep]
	\item 1 kg Mehl (Type 00 mit mindestens 12\% Eiweiß)
	\item 650 g kaltes Leitungswasser (Anfänger 600 g, weil leichter zu 
	verarbeiten)
	\item 25 g feines Salz
	\item 15 g Olivenöl
	\item 2 g frische Hefe
\end{itemize} 

\paragraph{Zutaten Pizzasoße}
 
\begin{itemize}[noitemsep]
	\item 500 g frische Tomaten
	\item 500 g passierte Tomaten aus der Flasche (hier von Mutti)
	\item 2 große Zwiebeln
	\item 2 Stück chinesischer Knoblauch
	\item 50 ml Olivenöl
	\item Salz
	\item je nach Geschmack Kräuter
\end{itemize}

\paragraph{Zubereitung }
Das Mehl in einer Rührschüssel vorlegen. Von dem Wasser ca. 100 ml 
abnehmen um die Hefe aufzulösen. Den Rest des Wassers zu dem Mehl 
geben 
und das Ganze mit einem Kochlöffel verrühren bis das gesamte Wasser 
aufgenommen wurde. Den Teig  abdecken und 30 Minuten ruhen lassen. Nach 
dem Ruhen das Olivenöl und das Wasser mit der aufgelösten Hefe zugeben und 
alles kneten, bis das Wasser und das Öl aufgenommen wurde, nun das Salz 
zugeben und den Teig fertig kneten ca. 15 Minuten. Da kann natürlich auch in 
einer Küchenmaschine erledigt werden. Knetzeit ca. 6 Minuten bei mittlerer 
Geschwindigkeit.
Den Teig in der Rührschüssel belassen und diese mit Frischhaltefolie 
verschließen. Den Teig 24 Stunden im Kühlschrank gehen lassen. Den Teig 
nach 24 Stunden aus der Schüssel nehmen, kurz durchkneten und in 6 
gleichmäßige Teiglinge teilen. Die Teiglinge zu Kugeln formen, in eine 
Auflaufform legen, diese mit Frischhaltefolie verschließen und nochmals 24 
Stunden im Kühlschrank gehen lassen. 

\paragraph{Zubereitung Tomatensoße}
Die Zwiebeln klein schneiden und im Olivenöl andünsten, nach einiger Zeit den 
gehackten Knoblauch zugeben. In der Zwischenzeit die frischen Tomaten in 
Stücke schneiden und zusammen mit den passierten Tomaten zu den 
angedünsteten Zwiebeln geben und köcheln lassen bis die Tomaten verkocht 
sind. Je nach Geschmack die Soße stückig lassen oder durch ein Sieb streichen.
Falls gewünscht Kräuter zugeben und mit Salz abschmecken.

\paragraph{Zubereitung Pizza}
Die Teiglinge aus der Form nehmen und von der Mitte heraus flach drücken, 
darauf achten, dass die Luft nicht aus den Rand gedrückt wird. Wenn der Boden 
die richtige Größe hat mit Tomatensoße bestreichen, nach Belieben belegen 
und bei mindestens 300°C im Holzbackofen backen.
\newpage

\begin{figure}[htbp]
	\centering
	\begin{minipage}{1\textwidth}
		\centering
		\includegraphics[width=1\linewidth]{pics/Pizza}
		\captionof{figure}{Pizza aus dem Holzbackofen}
		\label{fig:Pizza}
	\end{minipage}
\end{figure}
\newpage


\subsection{Brioche Buns}\label{BriocheBun}
Die Kunst ist es ein Bun herzustellen, der nicht zäh sondern schön mürbe ist. Er darf aber nicht zu schnell durchsuppen und sollte nicht in der 
Hand zerbröseln. Diese Eigenschaften in einem Brioche Buns zu vereinen kostet einiges an Experimentierarbeit. Zumal der Bun auch nicht 
trocken oder strohig werden darf.
Ich habe hier mein Rezept an ein Rezept aus dem Buch "Burger unser" (erschienen im Callwey Verlag, hier der Link zur in Leder gebunden Limited Edition \url{https://www.callwey.de/buecher/burger-unser-limited-edition/}) angelehnt, das für mich als gelernten Bäcker dem perfekten Brioche Bun  wohl am nächsten kommt.

\paragraph{Geräte}

\begin{itemize}[noitemsep]
	\item Holzbackofen
	\item Kenwood Cooking Chef
\end{itemize}

\paragraph{Zutaten}

\begin{itemize}[noitemsep]
	\item 3 EL Milch (3,8\% Fett)
	\item 50 g Rohrzucker
	\item 380 g Mehl, Type 550
	\item 1/2 Würfel frische Hefe
	\item 1 großes Ei
	\item 35 g warme Butter
	\item 45 g Mehl, Type 405
	\item 9 g Himalaya-Salz
	\item 240 ml Wasser
\end{itemize}

\paragraph{Zubereitung}

Milch auf ca. erwärmen nicht erhitzen, in eine Schüssel geben 1 EL Zucker und 1 EL Mehl
zugeben. Die Hefe hinein bröseln und alles verrühren. Stehen lassen
bis sich Blasen bilden. 
Die weiche Butter mit beiden Mehlsorten und Salz mischen bis eine streuselartige
Masse entsteht. Anschließend die Hefe-Milch-Mischung, Wasser und das Ei
langsam ein arbeiten. Die Küchenmaschine auf die mittlere Stufe einstellen und 
den Teig für ca. 8 Minuten kneten. Den Teig abdecken und 1-2 Stunden stehen lassen
bis sich das Volumen verdoppelt hat (Stockgare).
Nach der Stockgare den Teig auf der leicht bemehlten Arbeitsfläche in 8 gleiche Stücke
teilen. Die 8 Teile in schöne Kugeln formen, auf ein Backblech mit genügend Abstand legen
und eine Stunde liegen lassen, bis das Volumen verdoppelt hat (Stückgare).

Wenn die Buns aufgegangen sind mit einer Mischung aus 1 Eiklar, 1 EL Wasser und einer 
Prise Salz einstreichen. Dazu einen Silikonpinsel benutzen. Die Buns mit einem beliebigen
Topping, ich bevorzuge Sesam mit seinem unvergleichlichem nussigen Aroma, bestreuen
und 15-20 Minuten bei 200°C im Holzbackofen backen, bis die Oberfläche goldbraun ist.

\subsection{Kartoffel-Bun}
Der Kartoffel-Bun ist in Hinsicht der Luftigkeit und Stabilität die beste Alternative. Dieser 
Bun wird in den USA und in Kanada als Standard eingesetzt. Dieser Teig ist nicht so leicht 
wie der Brioche-Teig zu handhaben, ist allerdings deutlich luftiger. Es schmeckt auch 
weniger süßlich und ist für saftige Burger die erste Wahl. Auch dieses Rezept habe ich
an ein Rezept aus der selben Quelle wie das des Brioche Buns angelehnt.

\paragraph{Geräte}

\begin{itemize}[noitemsep]
	\item Holzbackofen
	\item Kenwood Cooking Chef
	\item Kartoffelstampfer oder Spätzlepresse
\end{itemize}

\paragraph{Zutaten}

\begin{itemize}[noitemsep]
	\item 300 ml Milch (3,8\% Fett)
	\item 1/2 Würfel frische Hefe
	\item 70 g brauner Zucker
	\item 300g Mehl, Type 405
	\item 200 g gekochte mehlige Kartoffeln (vom Vortag)
	\item 250 g Mehl, Type 550
	\item 10 g Himalaya-Salz
	\item 1 großes Ei
	\item 80 g Butter, weich
\end{itemize}

\paragraph{Zubereitung}

Die Milch auf erwärmen nicht erhitzen, in eine Schüssel geben 1 EL Zucker, 1 EL Mehl
und die Hefe zugeben und alles verrühren. Stehen lassen bis sich Blasen bilden. 

Die gekochten und geschälten Kartoffeln stampfen, bis keine groben 
Stücke mehr vorliegen oder durch eine Spätzlepresse drücken, aber niemals mit einem 
Handmixer verrühren, weil sich dadurch langkettige Moleküle bilden und die Kartoffeln
 Fäden ziehen. Die Kartoffeln zusammen mit der Hefe-Milch-Mischung, 
dem gesamten Mehl, dem Zucker, dem Salz und dem Ei in eine Küchenmaschine geben. 
Auf langsamer Stufe kneten bis ein grober Teig entsteht. 

Jetzt die Butter nach und nach einarbeiten, dabei die Küchenmaschine auf mittlere Stufe 
stellen. Anschließend die Küchenmaschine auf die höchste Stufe  einstellen und den 
Teig für 10 Minuten kneten. Den Teig in der Schüssel belassen, abdecken und für 1 bis 1,5
Stunden stehen lassen. Das Volumen sollte sich bis dahin verdoppelt haben (Stockgare).

Den Teig aus der Schüssel nehmen und gut kneten, dass die Luft entweicht. Den Teig 
mit den Händen flachdrücken und einschlagen, den Vorgang ein bis zweimal wiederholen.
Danach in 12 gleichmäßige Teile teilen und zu schönen kompakten Kugeln formen.

Die 12 Kugeln auf ein Backblech mit mit genügend Abstand legen und eine Stunde liegen
lassen, bis sich das Volumen verdoppelt hat (Stückgare).

Wenn die Buns aufgegangen sind mit einer Mischung aus 1 Eiklar, 1 EL Wasser und einer 
Prise Salz einstreichen. Dazu einen Silikonpinsel benutzen. Die Buns mit einem beliebigen
Topping, Schwarzkümmel mit ein wenig Salz ist auch ein schmackhaftes Topping mit einer
tollen Aromatik und 15-20 Minuten bei 200°C im Holzbackofen backen, bis die 
Oberfläche goldbraun ist.
